\subsection{Nearest work: retrieval-grounded FCA and closure-based knowledge expansion}

A particularly close 2026 result makes the novelty boundary around saturation and closure much sharper. Yang and Lee combine retrieval with Formal Concept Analysis (FCA) to construct a growing formal context in which a small language-model oracle evaluates incidences and proposed implications, returns counterexamples, and makes accepted implications, contradictions and corrections inspectable \cite{yanglee2026fca}. Their loop is explicitly verification-oriented rather than unrestricted text generation. Within its controlled object--attribute world, FCA supplies closure-based implication structure, and the experimental process expands the formal context across repeated rounds.

That work owns several functions that must not be redescribed as RAKL inventions: retrieval-grounded symbolic knowledge expansion, counterexample-driven correction, closure-based implication checking, and inspectable growth of a formal context. It also connects the present paper directly to the established FCA lineage, in which Galois connections and closure operators generate concept lattices \cite{ganter1999}. Earlier lattice-based knowledge-representation work used Birkhoff's representation theorem to implement finite distributive knowledge structures incrementally \cite{oles2000}. Likewise, constraint-processing research has long distinguished local consistency from conditions sufficient for global consistency \cite{dechter1992}. The parity obstruction in Section~4 should therefore be read as an explicit compact countermodel for RAKL's typing discipline, not as the discovery that local consistency can fail to become global consistency.

The difference in scope is structural. Yang and Lee study a bounded FCA formal context with a controlled attribute set and retrieval-grounded oracle. Epistemic mechanics treats the formal context itself as only one possible scientific representation and adds governance dimensions that FCA does not by itself supply: evidence-root independence, multi-coordinate authority, context/alignment transitions, immutable negative history, proposal-only workspace access, open-world route diversification, provenance-bearing omission and nearest-work audits, unresolved research fibers, and freshness-expiring saturation certificates. Conversely, RAKL does not currently implement the complete FCA attribute-exploration machinery or claim that its compatibility complex is a canonical implication basis.

This comparison also prevents a misleading use of the word ``saturated.'' FCA closure can be exact relative to a fixed formal context. The RAKL paper-level certificate is intentionally weaker and broader: it asks whether all registered substantive-growth coordinates have remained flat across heterogeneous expansion attempts under a fixed basis. A new source or mechanism can enlarge the effective research world and revoke that certificate even when the previously constructed FCA closure remains mathematically correct for its original context. Thus the nearest-work relation is complementary rather than competitive: FCA provides a rigorous closure engine inside a bounded representation, while epistemic saturation governs when a broader evidence-bearing research state is stable enough to release.
